\DocumentMetadata{
  pdfversion=2.0,
  pdfstandard=ua-2,
  lang=en-US,
  tagging=on,
  tagging-setup={math/setup=mathml-SE,tabsorder=structure}
}
\documentclass[11pt,letterpaper]{article}
\usepackage[margin=0.68in,includefoot]{geometry}
\usepackage{newcomputermodern}
\usepackage{amsmath,amscd,mathtools}
\usepackage{tikz,adjustbox}
\providecommand{\square}{\mdlgwhtsquare}
\usepackage[hidelinks]{hyperref}
\hypersetup{
  pdftitle={Logic PhD Exam, September 2004},
  pdfauthor={University of Florida Department of Mathematics},
  pdfsubject={Graduate mathematics examination},
  pdfdisplaydoctitle=true,
  bookmarksopen=true
}
\setcounter{secnumdepth}{0}
\setlength{\parindent}{0pt}
\setlength{\parskip}{0.28em}
\setlength{\emergencystretch}{3em}
\setlength{\leftmargini}{2.15em}
\setlength{\leftmarginii}{2.65em}
\setlength{\leftmarginiii}{3em}
\renewcommand{\labelenumi}{\textbf{\theenumi.}}
\pagestyle{empty}
\raggedbottom
\allowdisplaybreaks
\newenvironment{examproblems}[1]{%
  \begin{enumerate}%
  \setcounter{enumi}{#1}%
  \setlength{\topsep}{0.35em}%
  \setlength{\partopsep}{0pt}%
  \setlength{\itemsep}{0.48em}%
  \setlength{\parsep}{0.18em}%
}{\end{enumerate}}
\newenvironment{examsubparts}{%
  \begin{enumerate}%
  \setlength{\topsep}{0.18em}%
  \setlength{\partopsep}{0pt}%
  \setlength{\itemsep}{0.16em}%
  \setlength{\parsep}{0.1em}%
}{\end{enumerate}}
\newenvironment{examinstructions}{%
  \par\begingroup\itshape
}{%
  \par\endgroup\vspace{0.18em}
}
\makeatletter
\renewcommand\section{\@startsection{section}{1}{0pt}%
  {-1.25ex plus -.25ex minus -.1ex}%
  {0.55ex plus .1ex}%
  {\normalfont\large\bfseries}}
\renewcommand{\@maketitle}{%
  \newpage\null\vskip -1em%
  {\footnotesize\bfseries
    \href{https://www.ufl.edu/}{University of Florida}, \href{https://math.ufl.edu/}{Department of Mathematics}\par}%
  \vskip -0.5em%
  \tagstructbegin{tag=Title}%
    {\LARGE\bfseries\href{https://gma.math.ufl.edu/exams/logic-phd/}{Logic PhD Exam}, September 2004\par}%
  \tagstructend%
  \vskip 0.25em%
  \tagmcbegin{artifact}%
    \hrule height 1pt%
  \tagmcend%
  \vskip 1em%
}
\makeatother
\title{Logic PhD Exam, September 2004}
\author{}
\date{}
\begin{document}
\maketitle
\thispagestyle{empty}
\begin{examinstructions}
Answer 7 questions.
\end{examinstructions}
\begin{examproblems}{0}
\item
Sketch a proof that every truth function \(E:\{T,F\}^n\to\{T,F\}\) can be represented by a propositional sentence.
\item
State the (Countable) Compactness for First Order Logic and sketch the proof using Henkin completeness.
\item
Sketch a proof that the theory of dense linear orderings without end points satisfies Quantifier Elimination and hence is complete and decidable.
\item
State and prove the Omitting Types Theorem.
\item
Show that Zorn’s Lemma implies the Axiom of Choice.
\item
Show that for any notion of forcing~\(P\) and any countable set~\(D\) of~\(P\)-dense sets, there exists a~\(D\)-generic~\(P\)-filter.
\item
Define the constructible universe~\(L\) and sketch a proof that~\(L\) satisfies the Continuum Hypothesis.
\item
Sketch a proof of the Undecidability of (Peano) Arithmetic.
\item
Show that an infinite set \(X\subseteq\omega\) is the domain of a partial recursive function if and only if~\(X\) is semirecursive (that is, \(x\in X\Longleftrightarrow(\exists y)R(x,y)\) where~\(R\) is a recursive relation).
\item
Show how to construct a simple recursively enumerable set which is neither recursive nor~\(m\)-complete.
\end{examproblems}
\end{document}
